
\def\arraystretch{1.3}

\begin{rubric}{Current Research Projects}

\entry*[] ``PLM-based Sentiment Extraction and Reference Letters.'' Arxiv Working Paper \url{https://arxiv.org/abs/2410.16325}.

\entry*[] ``Measuring Product Quality Using Images: The CLIP Q-Score and an Application to Real Estate.'' Arxiv Working Paper \url{https://arxiv.org/abs/2608.01544}.

%  {\footnotesize
%      \emph{Abstract}: We use transformer models that have been pre-trained to interpret images as well as text to extract a \emph{luxury score} for a large number of real estate properties. First, we developed a web-crawling technique to keep track of the largest real estate platform in Russia. As a result, we have collected a data set consisting of almost the entirety of properties that were offered for sale or rent in Moscow during the summer of 2025. Besides all the usual property descriptors, these data contains tens of thousands of user-posted photos of the properties. We created a completely new approach to the extraction of information from images. We take advantage of the fact that CLIP (a multimodal large language model) was pre-trained using matched image-caption pairs taken from the internet. Through its training, the model's main task is to learn to distinguish correct and incorrect captions for the image. Our approach exploits this training by submitting each image in our data to CLIP, together with different text prompts, each of which describes the property as increasingly luxurious. We then combine the model's outputs in a summary measure of the state of repair/disrepair of the property. Our luxury score follows geographic and other patterns completely compatible with its interpretation as a measure of real estate quality: 1) The best real estate is in the center, north-western and western administrative districts; the properties is the far east and far south districts is in terrible state of disrepair. 2) Properties situated in buildings scheduled for demolition are in significantly worse condition than neighboring properties. 3) There is a U-shaped pattern of real estate quality relative of year of construction, consistent with the historical record regarding housing investment during the Soviet era. 4) Owner self-reported `state of repair' is highly positively correlated with our luxury score.
%  }

\entry*[] ``Good Looks and the Academic Job Market.''

\entry*[] ``A Prompt-based Approach to Manager Sentiment.''

% \begin{quotation}
%      \emph{Abstract}: The idea that agents' sentiment (understood as a subjective state of optimism/pesimism regarding the future of the economy) is an important determinant of investment has a long tradition in macro-finance. We apply a prompt-based sentiment extraction technique (Slonimczyk 2025) to obtain a \emph{macro measure of sentiment} from regulatory filings by publicly traded companies (10-K and 10-Q).
% \end{quotation}


% \begin{quotation}
%      \emph{Abstract}:  There is extensive and convincing evidence linking individual appearance to income, promotion in the career ladder, and self-reported happiness. One drawback, however, is that the literature widely relies on human judgement to assess beauty. We develop an innovative approach to measure ``good looks'' in a way that is relatively objective and fully reproducible. Our approach involves the use of visual transformers (ViTs), i.e. large language models that have been adapted to processing images. Such models are able to encode images, effectively mapping them into relatively low-dimensional numerical spaces. We use different ViTs to encode the images of job candidates in the economics academic market. Key to our approach to measuring beauty is that we also encode the images of a large number of celebrities. Our definition of good looks is simply the average distance between the embeddings of the candidate's photos and the embeddings corresponding to celebrities. We validate our approach using the Chicago Face Dataset, a well-know repository of portrait-style photos that were scored by 26 independent human raters. Our beauty scores not only are closely correlated with human ratings, but are also statistically significant predictors of job market success. More notably, the beauty score is highly positively correlated with the sentiment score of the letters of recommendation, suggesting that one of the channels for good looks is through access to connections.
%    \end{quotation}

\entry*[] ``Differential Treatment of Families with Children.''

% \begin{quotation}
%      \emph{Abstract}:  We conduct a series of field experiments designed to test whether there is discrimination against children in the real estate market. Phase 1 has received approval from HSE's Committee on Research Ethics and has been pre-registered in the AEA RCT Registry. Phase 2 is scheduled for the Fall 2026.
% \end{quotation}

\end{rubric}
